\chapter{Weakly étale subalgebras of fields}
\label{ch:mathlib-weakly-etale-subalgebra}

This chapter backs the infrastructure file
\texttt{Proetale/Mathlib/RingTheory/WeaklyEtale/Subalgebra.lean}. It hosts a
field-base specialisation of the Stacks 092Q / 0CKR(3) statement that the
consumer \texttt{Algebra.WeaklyEtale.isAlgebraic\_of\_isField\_aux} in
\texttt{FieldExtension.lean} requires to close the transcendental case of
Stacks 092P.

\section{Statement}

\begin{lemma}
  \label{lem:weakly-etale-adjoin-singleton}
  \lean{Algebra.WeaklyEtale.adjoin_singleton}
  Let $K$ be a field and $L$ a $K$-algebra with $K \to L$ weakly étale.
  For every $a \in L$, the $K$-subalgebra $K[a] = \mathrm{adjoin}_K\{a\}$
  is weakly étale over $K$.
  \uses{def:weakly-etale-algebra}
\end{lemma}

Recall that \emph{weakly étale} means $K$-flat together with formally
unramified. The flatness half is automatic: $K$ is a field, so every
$K$-module is free and hence flat, and in particular every $K$-subalgebra
of $L$ is $K$-flat. The substantive content is formal unramifiedness; the
present chapter discharges it by reducing to the stronger property of
étaleness for finitely generated $K$-subalgebras of a weakly étale
$K$-algebra.

\section{Strategy of proof}

The actual closure used in the formalisation is a short consequence chain
that delegates the substantive content to two prior results in the
project:

\begin{enumerate}
  \item The subalgebra $K[a] = \mathrm{adjoin}_K\{a\}$ is finitely
    generated over $K$, being generated by the single element $a$.
  \item By
    \Cref{thm:weakly-etale-imples-ind-etale-over-fields} (Stacks 092Q,
    item (1); the underlying statement is also recorded as Stacks
    tag 0CKR(3) / Stacks 15.105.8), every finitely generated
    $K$-subalgebra of a weakly étale $K$-algebra is étale over $K$.
    Applying this to $A := K[a]$ yields that $K \to A$ is étale.
  \item Étale algebras are weakly étale
    (\Cref{lemma:etale-weakly-etale-algebra}), which gives the claim.
\end{enumerate}

The substantive Stacks 092Q content (every FG $K$-subalgebra of a
weakly étale $K$-algebra is étale) is proved in a separate chapter and
its Lean home is
\texttt{Proetale/Algebra/WeaklyEtaleField.lean::etale\_of\_fg}. The
present chapter therefore only carries the short reduction; the
descent-of-idempotency / absolute-flatness machinery cited there is not
re-derived here.

\paragraph{Alternative route (not used).}
A more hands-on proof avoids the FG-subalgebra theorem by descending
idempotency of $\ker \mu_L$ along the $K$-flat injection
$K[a] \otimes_K K[a] \hookrightarrow L \otimes_K L$, using
$\mathrm{KaehlerDifferential.ideal\_eq\_diagonal}$ to expose
$\ker \mu_{K[a]}$ as the ideal generated by the diagonal elements
$\{1 \otimes b - b \otimes 1 : b \in K[a]\}$ and then transporting the
$(1-e)$-annihilation property from the overring. This route was the
original Lane B target in earlier iterations; it is recorded here for
future readers but is not the route taken by the current Lean proof.

\section{Localization is weakly étale}

The second Mathlib-gap statement carried in this file is the very basic
fact that any localization $R \to S^{-1} R$ at a multiplicative submonoid
is weakly étale. It is one of the three ``standard sources of weakly
étale maps'' recorded by Stacks 092J (the other two being étale maps and
filtered colimits of weakly étale maps), and the project needs it as a
named lemma in order to feed the formal-étale / pro-étale infrastructure
that consumes weakly étale ring maps in bulk.

\begin{lemma}\leanok
[\stacksproject{092J}, localization clause]
  \label{lem:weaklyEtale-localization}
  \lean{Algebra.WeaklyEtale.localization}
  \uses{def:weakly-etale-algebra}
  % SOURCE: Stacks Project, Tag 092J, item (1) (label `lemma-when-weakly-etale'
  % inside Section 15.106 "Weakly étale ring maps" of `more-algebra.tex')
  % (read from references/stacks-092PQ.tex, lines 290--299).
  % SOURCE QUOTE: "Let $A \to B$ be a ring map. Then $A \to B$ is weakly \'etale in each
  % of the following cases
  % \begin{enumerate}
  % \item $B = S^{-1}A$ is a localization of $A$,
  % \item $A \to B$ is \'etale,
  % \item $B$ is a filtered colimit of weakly \'etale $A$-algebras.
  % \end{enumerate}"
  \textit{Source: Stacks Project, Tag 092J (item (1)).}
  Let $R$ be a commutative ring and $S \subseteq R$ a multiplicative
  submonoid. Then the localization map $R \to S^{-1} R$ is weakly étale.
\end{lemma}

% SOURCE QUOTE PROOF: "An \'etale ring map is flat and the map
% $B \otimes_A B \to B$ is also \'etale as a map between \'etale
% $A$-algebras (Algebra, Lemma \ref{algebra-lemma-map-between-etale}).
% This proves (2).
%
% \medskip\noindent
% Let $B_i$ be a directed system of weakly \'etale $A$-algebras.
% Then $B = \colim B_i$ is flat over $A$ by Algebra, Lemma
% \ref{algebra-lemma-colimit-flat}. Note that the transition maps
% $B_i \to B_{i'}$ are flat by Lemma \ref{lemma-weakly-etale-permanence}.
% Hence $B$ is flat over $B_i$ for each $i$, and we see that $B$ is flat
% over $B_i \otimes_A B_i$ by Algebra, Lemma
% \ref{algebra-lemma-composition-flat}. Thus $B$ is flat over
% $B \otimes_A B = \colim B_i \otimes_A B_i$ by Algebra, Lemma
% \ref{algebra-lemma-colimit-rings-flat}.
%
% \medskip\noindent
% Part (1) can be proved directly, but also follows by combining (2) and (3)."
\begin{proof}
  \uses{def:weakly-etale-algebra}
  \leanok
  \textit{Source: Stacks Project, Tag 092J (item (1)).}

  We unwind Mathlib's definition of \emph{weakly étale}. The class
  \texttt{Algebra.WeaklyEtale}~$R\,B$ packages two flatness conditions
  on the structure map $R \to B$:
  \begin{enumerate}
    \item[(F1)] $B$ is flat as an $R$-module;
    \item[(F2)] $B \otimes_R B$ is flat as a $B$-module via the
      multiplication map $\mu_B : B \otimes_R B \to B$ (equivalently
      via either factor inclusion, which agree up to the diagonal).
  \end{enumerate}
  We verify (F1) and (F2) for the localization $B := S^{-1} R$.

  \emph{Step 1 (verification of (F1)).} The structure map
  $R \to S^{-1} R$ is flat. This is the standard property of
  localization-as-base-change: $S^{-1} R$ is the localization of $R$ at
  the multiplicative set $S$, hence a flat $R$-module. In Mathlib this
  is recorded as \texttt{IsLocalization.flat} (for any
  \texttt{IsLocalization $S$ $T$} instance, $T$ is $R$-flat).

  \emph{Step 2 (verification of (F2)).} We must show that
  $(S^{-1} R) \otimes_R (S^{-1} R)$ is flat over $S^{-1} R$. Localization
  tensors are again localizations: there is a canonical $S^{-1} R$-algebra
  isomorphism
  \[
    (S^{-1} R) \otimes_R (S^{-1} R) \;\cong\; S^{-1}(S^{-1} R),
  \]
  where on the right we view $S^{-1} R$ as an $R$-algebra and localize
  again at the image of $S$. (In fact $S^{-1}(S^{-1} R) \cong S^{-1} R$
  because $S$ already maps to units in $S^{-1} R$, but we do not need this
  refinement — the bare localization identity suffices.) The right-hand
  side is, by construction, the localization of $S^{-1} R$ at the
  multiplicative submonoid generated by the image of $S$ in $S^{-1} R$.
  Applying \texttt{IsLocalization.flat} again — this time with
  \emph{base ring $S^{-1} R$} — yields that
  $S^{-1}(S^{-1} R)$, and hence $(S^{-1} R) \otimes_R (S^{-1} R)$, is
  flat as an $S^{-1} R$-module.

  \emph{Step 3 (assembly).} Conditions (F1) and (F2) together unfold
  the Mathlib definition of \texttt{Algebra.WeaklyEtale}~$R\,(S^{-1} R)$.
  Hence $R \to S^{-1} R$ is weakly étale, which is the claim.

  Stacks 092J derives item (1) by combining items (2) and (3): a
  localization is an étale map composed with itself in the trivial
  filtered-colimit sense. The proof recipe above (direct verification
  of the two Mathlib flatness conditions) is the version recorded as
  ``Part (1) can be proved directly'' in the verbatim Stacks proof
  above; it is what the Lean formalisation in
  \texttt{Proetale/Mathlib/RingTheory/WeaklyEtale/Localization.lean}
  carries out.
\end{proof}

\section{Downstream consumer (informal)}

The bivariate contradiction inside
\texttt{isAlgebraic\_of\_isField\_aux} consumes
\Cref{lem:weakly-etale-adjoin-singleton} as follows. Suppose $a \in L$
is transcendental over $K$; then $K[a] \cong K[X]$ and
$K[a] \otimes_K K[a] \cong K[X, Y]$, which is a domain. By
\Cref{lem:weakly-etale-adjoin-singleton}, $K \to K[a]$ is weakly étale,
so $\mu_{K[a]}$ is formally étale and $\ker \mu_{K[a]}$ is principal,
generated by an idempotent of $K[X, Y]$. The only idempotents of a
domain are $0$ and $1$, so $\ker \mu_{K[a]} = (0)$ or $(1)$. But
$\ker \mu_{K[a]} = (X - Y)$, which is neither — contradiction. Hence
$a$ is algebraic.
